\documentclass[11pt,a4paper]{article}

\usepackage[utf8]{inputenc}
\usepackage[T1]{fontenc}
\usepackage[french]{babel}
\usepackage{amsmath,amssymb,amsfonts}
\usepackage{geometry}
\usepackage{graphicx}
\usepackage{booktabs}
\usepackage{hyperref}
\usepackage{array}
\usepackage{xcolor}
\usepackage{float}

\geometry{margin=2.5cm}
\graphicspath{{figures/}}

\title{
\textbf{Des graphes d'intrication aux équations d'Einstein}\\
\large Dérivation effective de la gravité continue dans Bottom-Up Quantum Gravity
}

\author{Farid Hamdad}
\date{\today}

\begin{document}

\maketitle

\begin{abstract}
Dans les papiers précédents du programme Bottom-Up Quantum Gravity, la
géométrie émergente a été reconstruite à partir de la matrice d'information
mutuelle \(W_{ij}=I(i:j)\), via le Laplacien d'intrication \(L_{\rm ent}\).
Paper~14 a montré que le terme spectral de l'action BuP n'est plus libre :
l'exposant \(\beta\) est contraint par les invariants spectraux du graphe.
Le présent papier étudie l'étape suivante : les conditions sous lesquelles
l'équation variationnelle discrète
\[
\frac{\delta S_{\rm BuP}[W]}{\delta W_{ij}}=0
\]
admet une limite continue de type Einstein,
\[
G_{\mu\nu}[g^{\rm ent}]
+
\Lambda_{\rm ent}g_{\mu\nu}^{\rm ent}
=
8\pi G_{\rm eff}T_{\mu\nu}^{\rm ent}.
\]
Nous montrons numériquement trois briques nécessaires à cette limite. Le
Laplacien rescalé \(L_\epsilon=(D-W)/\epsilon\) récupère la dimension
spectrale attendue sur des géométries contrôlées, avec
\(d_s({\rm circle})=1.0021\) et \(d_s({\rm grid2d})=1.9938\). La courbure
d'Ollivier--Ricci distingue une sphère d'un tore plat périodique, avec un
excès relatif \(\Delta\bar\kappa_{\rm sphere-flat}=0.010525\). Enfin, une
perturbation radiale locale de \(W_{ij}\) produit une réponse de courbure
localisée, avec des ratios near/far de \(7.75\) sur le tore plat et \(11.51\)
sur la sphère. Ces résultats ne constituent pas encore une preuve analytique
des équations d'Einstein, mais fournissent une validation numérique contrôlée
des trois flèches nécessaires à la dérivation effective :
\[
W_{ij}\to L_\epsilon\to\Delta_g,
\qquad
\kappa_{ij}^{\rm OR}\to R_{\mu\nu}u^\mu u^\nu,
\qquad
\delta W_{\rm loc}\to\delta\kappa(r).
\]
\end{abstract}

\tableofcontents

\section{Introduction}

Le programme Bottom-Up Quantum Gravity repose sur l'idée que l'espace, le
temps et la gravité ne sont pas des structures fondamentales, mais des
régimes effectifs issus de la structure d'intrication d'un état quantique
global. L'objet fondamental n'est pas une métrique \(g_{\mu\nu}\), mais une
matrice relationnelle
\[
W_{ij}=I(i:j),
\]
où \(I(i:j)\) désigne l'information mutuelle entre les degrés de liberté
\(i\) et \(j\).

À partir de \(W_{ij}\), on construit le Laplacien d'intrication
\[
L_{\rm ent}=D-W,
\]
où
\[
D_{ii}=\sum_j W_{ij}.
\]
Ce Laplacien encode la connectivité, la diffusion, la dimension spectrale
et les propriétés géométriques émergentes du réseau d'intrication.

Les papiers précédents ont établi plusieurs résultats intermédiaires :
\begin{itemize}
    \item la dimension spectrale \(d_s\) et la dimension de marche \(d_w\)
    définissent un exposant gravitationnel effectif ;
    \item la courbure discrète d'Ollivier--Ricci réagit aux perturbations
    locales d'information ;
    \item Paper~8 a construit un tenseur énergie-impulsion effectif à partir
    de perturbations locales \(\delta W_{\rm loc}\) ;
    \item Paper~14 a montré que la dynamique modulaire et la gravité
    effective sont contrôlées par les mêmes invariants spectraux.
\end{itemize}

Le présent papier franchit une étape supplémentaire. Il pose la question :

\[
\boxed{
\text{Sous quelles conditions l'équation discrète BuP converge-t-elle vers
une équation d'Einstein effective ?}
}
\]

La cible est la chaîne suivante :
\[
\frac{\delta S_{\rm BuP}[W]}{\delta W_{ij}}=0
\quad
\xrightarrow[N\to\infty]{}
\quad
G_{\mu\nu}[g^{\rm ent}]
+
\Lambda_{\rm ent}g_{\mu\nu}^{\rm ent}
=
8\pi G_{\rm eff}T_{\mu\nu}^{\rm ent}.
\]

Il ne s'agit pas ici de postuler les équations d'Einstein, mais d'étudier
leur émergence comme limite continue d'une condition d'équilibre sur un réseau
d'intrication.

\section{Point de départ : action BuP discrète}

La formulation variationnelle minimale de BuP est :
\[
S_{\rm BuP}[W]
=
\underbrace{\mathrm{Tr}\,L(W)^{-\beta[W]}}_{T_1}
+
\underbrace{\sum_{ij}W_{ij}\kappa_{ij}[W]}_{T_2}
+
\underbrace{\lambda\sum_{ij}W_{ij}d_{ij}^{2}}_{T_3}
+
\underbrace{S_{\rm topo}[W]}_{T_4}.
\]

Chaque terme joue un rôle distinct.

\subsection{Terme spectral}

Le terme
\[
T_1=\mathrm{Tr}\,L(W)^{-\beta[W]}
\]
encode l'action géométrique globale du graphe d'intrication. Paper~14 a
montré que \(\beta\) n'est plus un paramètre libre, mais une fonction des
invariants spectraux :
\[
\beta
=
F(\lambda_2,d_s,d_w,\Delta_3).
\]

La relation effective obtenue dans Paper~14 est :
\[
\beta_{\rm mod}
=
(0.531\pm0.033)
+
(1.726\pm0.118)\,\beta_{\rm grav},
\]
avec
\[
\beta_{\rm grav}
=
\frac{\alpha_{\rm eff}+1}{2},
\qquad
\alpha_{\rm eff}
=
\frac{2d_s}{d_w}+d_w-4.
\]

Ainsi, le terme spectral se ferme sur le spectre du graphe lui-même.

\subsection{Terme de courbure}

Le terme
\[
T_2=\sum_{ij}W_{ij}\kappa_{ij}[W]
\]
encode la réponse locale de courbure du graphe, où \(\kappa_{ij}\) est la
courbure d'Ollivier--Ricci du lien \((i,j)\). Ce terme est l'ancêtre discret
du terme de courbure continue.

Dans une limite géométrique régulière, on s'attend à une correspondance du
type :
\[
\kappa_{ij}^{\rm OR}
\longrightarrow
R_{\mu\nu}u^\mu u^\nu,
\]
où \(u^\mu\) est la direction tangente associée au lien.

\subsection{Terme de localité}

Le terme
\[
T_3=\lambda\sum_{ij}W_{ij}d_{ij}^2
\]
pénalise les connexions longues et impose une structure locale. Sans ce terme,
le graphe optimal tendrait vers une connectivité trop dense, incompatible avec
une géométrie locale émergente.

\subsection{Terme topologique}

Le terme
\[
S_{\rm topo}[W]
\]
encode les contraintes globales de connectivité, de composantes, de cycles et
de topologie effective. Sa forme exacte reste ouverte. Il représente une des
directions de formalisation futures de BuP.

\section{Équation variationnelle discrète}

L'équation fondamentale est :
\[
\frac{\delta S_{\rm BuP}[W]}{\delta W_{ij}}=0.
\]

Plutôt que d'écrire explicitement une dérivée simplifiée du terme spectral,
on définit :
\[
\mathcal{G}_{ij}^{(\beta)}[L]
=
-\frac{\partial}{\partial W_{ij}}
\mathrm{Tr}\,L^{-\beta}.
\]

L'équation d'équilibre prend alors la forme :
\[
\mathcal{G}_{ij}^{(\beta)}[L]
=
\kappa_{ij}
+
\lambda d_{ij}^{2}
+
\frac{\partial S_{\rm topo}}{\partial W_{ij}}.
\]

Cette équation est l'analogue discret d'une équation géométrique : le
propagateur spectral du réseau d'intrication est équilibré par la courbure
locale, le coût de localité et les contraintes topologiques.

La question de Paper~15 est de savoir si, dans une limite continue stable,
cette équation converge vers :
\[
G_{\mu\nu}
+
\Lambda_{\rm ent}g_{\mu\nu}
=
8\pi G_{\rm eff}T_{\mu\nu}^{\rm ent}.
\]

\section{Distance d'intrication et métrique émergente}

Une distance informationnelle candidate est :
\[
d_{ij}^{\rm ent}
=
-\xi\log\left(
\frac{W_{ij}+\epsilon}{W_{\max}}
\right),
\]
où \(\epsilon\) régularise les liens nuls et \(\xi\) fixe l'échelle de
longueur.

Dans une limite lisse, on attend :
\[
d_{\rm ent}(x,y)^2
=
g_{\mu\nu}^{\rm ent}(x)
(x^\mu-y^\mu)(x^\nu-y^\nu)
+
O(|x-y|^4).
\]

La métrique peut alors être reconstruite localement à partir du développement
quadratique de la distance :
\[
g_{\mu\nu}^{\rm ent}(x)
\sim
\frac{1}{2}
\partial_{x^\mu}\partial_{y^\nu}
d_{\rm ent}(x,y)^2
\Big|_{y=x}.
\]

Numériquement, cette reconstruction peut être approchée par MDS, par
géométrie des distances locales, ou par ajustement local d'une forme
quadratique.

\section{Limite spectrale : \(L_N\to\Delta_g\)}

La première flèche nécessaire à la dérivation est :
\[
L_N\longrightarrow \Delta_g.
\]

Le Laplacien cible est le Laplacien de Laplace--Beltrami :
\[
\Delta_g
=
\frac{1}{\sqrt{|g|}}
\partial_\mu
\left(
\sqrt{|g|}g^{\mu\nu}\partial_\nu
\right).
\]

Le test numérique se fait via la trace de chaleur :
\[
Z(t)=\mathrm{Tr}\,e^{-tL},
\]
et la dimension spectrale :
\[
d_s(t)
=
-2\frac{d\log Z(t)}{d\log t}.
\]

Un premier test utilisant le Laplacien normalisé non rescalé
\[
L_{\rm norm}=I-D^{-1/2}WD^{-1/2}
\]
échoue sur les géométries 2D : les dimensions spectrales se replient vers
\(d_s\simeq1.2\). Ce résultat négatif montre que le graphe seul ne suffit
pas ; l'échelle de diffusion doit être prise en compte.

La version positive utilise le Laplacien rescalé :
\[
L_\epsilon=\frac{D-W}{\epsilon}.
\]

Avec \(k=\sqrt{N}\), \(\epsilon=0.5\,\epsilon_{\rm knn}\), et \(N=1024\), on
obtient :

\[
\begin{array}{lccc}
\toprule
{\rm Géométrie} & d_{\rm cible} & d_s^{\rm mesuré} & {\rm erreur} \\
\midrule
{\rm cercle} & 1 & 1.0021 & 0.0021 \\
{\rm intervalle} & 1 & 0.9445 & 0.0555 \\
{\rm grille\ 2D} & 2 & 1.9938 & 0.0062 \\
{\rm sphère} & 2 & 1.8964 & 0.1036 \\
\bottomrule
\end{array}
\]

Ces résultats fournissent une première évidence numérique positive pour :
\[
L_N\longrightarrow\Delta_g.
\]

\section{Courbure discrète et signal de Ricci}

La deuxième flèche nécessaire est :
\[
\kappa_{ij}^{\rm OR}
\longrightarrow
R_{\mu\nu}u^\mu u^\nu.
\]

La courbure d'Ollivier--Ricci d'un lien \((i,j)\) est définie par :
\[
\kappa_{ij}
=
1-
\frac{W_1(m_i,m_j)}{d(i,j)},
\]
où \(W_1\) est la distance de Wasserstein-1 entre les mesures de voisinage
\(m_i\) et \(m_j\).

Un premier test utilisant une grille plane avec bord produit une courbure
positive artificielle, due aux effets de frontière. La version v2 remplace la
grille par un tore plat périodique et utilise des distances géodésiques
intrinsèques.

À la résolution finale, le tore plat donne :
\[
\bar\kappa_{\rm flat}=-0.001671\simeq0,
\]
tandis que la sphère donne :
\[
\bar\kappa_{\rm sphere}=0.008854.
\]

Le signal relatif est donc :
\[
\Delta\bar\kappa_{\rm sphere-flat}=0.010525.
\]

Cela fournit une première indication numérique positive que la courbure
d'Ollivier--Ricci discrète détecte le signal de Ricci continu :
\[
\kappa_{ij}^{\rm OR}
\longrightarrow
R_{\mu\nu}u^\mu u^\nu.
\]

Ce résultat reste qualitatif : il ne constitue pas encore une preuve de
convergence ponctuelle. Il valide cependant la séparation entre géométrie
plate périodique et géométrie positivement courbée.

\section{Source effective : héritage de Paper 8}

Paper~8 a déjà construit une source effective depuis une perturbation locale
de l'intrication :
\[
\delta W_{\rm loc}
\longrightarrow
T_{\mu\nu}^{\rm eff}.
\]

La forme générale proposée était :
\[
T_{\mu\nu}^{\rm eff}
=
T_{\mu\nu}^{\rm matter}[\delta W^{\rm loc}]
+
T_{\mu\nu}^{\rm ent}[d_s].
\]

La meilleure source candidate dans Paper~8 était :
\[
S_{\rm flux}
=
T_{00}
-
\frac{1}{2}T_{aa}
+
\frac{1}{2}T_{\rm grad}
+
F.
\]

Elle présentait une corrélation de Spearman :
\[
\rho=0.741,
\qquad
p=1.84\times10^{-4},
\]
avec les fluctuations de courbure. Ce résultat constitue une preuve de concept
numérique pour la correspondance :
\[
\delta W_{\rm loc}
\longrightarrow
T_{\mu\nu}^{\rm eff}.
\]

Paper~8 montre également que le secteur matière n'est pas conservé isolément :
\[
\nabla^\mu T_{\mu\nu}^{\rm matter}\neq0.
\]

Dans BuP, cette non-conservation n'est pas interprétée comme un échec, mais
comme un échange avec le fond d'intrication :
\[
J_\nu^{\rm exchange}
=
\nabla^\mu
\left[
G_{\rm eff}(d_s)T_{\mu\nu}^{\rm matter}
\right],
\]
avec compensation :
\[
\nabla^\mu
\left[
G\,T_{\mu\nu}^{\rm ent}
\right]
=
-
J_\nu^{\rm exchange}.
\]

Le champ d'intrication utilisé comme proxy était :
\[
\phi_{\rm ent}
=
\delta d_s
+
\delta R
+
\frac{1}{2}\Delta\delta R.
\]

Ainsi, Paper~15 ne part pas d'une source inconnue. Il prend la construction
de Paper~8 comme point de départ et cherche à la formaliser dans une limite
continue.

\section{Réponse de courbure à une source d'intrication}

Le test source-response v1 modifiait uniquement les liens internes à la
région source :
\[
W'_{ij}=(1+s)W_{ij},
\qquad i,j\in{\rm source}.
\]
Il produisait une réponse ultra-locale.

La version v2 introduit une source radiale lisse :
\[
\phi_i
=
\exp\left[
-\frac{d(i,\mathrm{source})^2}{2\sigma^2}
\right],
\]
et perturbe les poids selon :
\[
W'_{ij}
=
W_{ij}
\left[
1+s\frac{\phi_i+\phi_j}{2}
\right].
\]

La réponse de courbure est :
\[
\Delta\kappa_{ij}
=
\kappa_{ij}^{\rm after}
-
\kappa_{ij}^{\rm before}.
\]

Pour \(s=-0.30\), on obtient sur le tore plat :
\[
\frac{
\langle|\Delta\kappa|\rangle_{\rm near}
}{
\langle|\Delta\kappa|\rangle_{\rm far}
}
=
7.75.
\]

Sur la sphère :
\[
\frac{
\langle|\Delta\kappa|\rangle_{\rm near}
}{
\langle|\Delta\kappa|\rangle_{\rm far}
}
=
11.51.
\]

La réponse suit fortement le profil de source. Pour \(s=-0.30\), la
corrélation entre \(\phi_{\rm edge}\) et \(|\Delta\kappa|\) vaut :
\[
\rho_{\rm Spearman}=0.583,
\qquad
p=4.12\times10^{-83}
\]
sur le tore plat, et :
\[
\rho_{\rm Spearman}=0.752,
\qquad
p=5.49\times10^{-165}
\]
sur la sphère.

De plus, la réponse signée change de signe lorsque \(s\) passe de négatif à
positif. Cela montre que la source d'intrication induit une réponse de
courbure signée et localisée.

Ce test fournit donc une réalisation numérique contrôlée de la chaîne :
\[
\delta W_{\rm loc}
\longrightarrow
\delta\kappa(r).
\]

\subsection{Convergence avec les résultats de Paper 7 et Paper 8}

Le résultat central du Step~E est que la réponse de courbure est à la fois
locale et corrélée au profil de perturbation d'intrication. Sur la sphère,
pour une perturbation \(s=-0.30\), on obtient :
\[
\frac{
\langle|\Delta\kappa|\rangle_{\rm near}
}{
\langle|\Delta\kappa|\rangle_{\rm far}
}
=
11.51,
\qquad
\rho_{\rm Spearman}
\left(
\phi_{\rm edge},
|\Delta\kappa|
\right)
=
0.752.
\]
Ce résultat montre que la perturbation locale de \(W_{ij}\) induit une réponse
de courbure fortement localisée, et que l'amplitude de cette réponse suit le
profil radial de la source d'intrication.

Ce signal rejoint deux résultats antérieurs du programme BuP. Paper~7 avait
montré, sur des graphes d'information mutuelle quantiques à \(N=16\), qu'une
perturbation locale de \(W_{ij}\) produit une réponse positive de courbure,
avec
\[
\Delta\kappa_{\rm edge}=0.076
\]
et une fraction positive de \(100\%\) dans le régime testé. Paper~8 avait
ensuite reconstruit un tenseur énergie--impulsion effectif depuis les
perturbations locales d'intrication, avec une corrélation de Spearman
\[
\rho=0.741
\]
entre la source reconstruite et les fluctuations de courbure.

Les trois mesures peuvent être résumées ainsi :
\[
\begin{array}{llll}
\toprule
{\rm Paper} & {\rm Test} & {\rm Cadre} & {\rm Signal}\ \delta W\to\delta\kappa \\
\midrule
{\rm Paper\ 7}
& {\rm Réponse\ de\ courbure}
& {\rm graphes\ MI\ quantiques}\ (N=16)
& \Delta\kappa_{\rm edge}=0.076,\ {\rm fraction\ positive}=100\% \\
{\rm Paper\ 8}
& T_{\mu\nu}^{\rm eff}\ {\rm reconstruit}
& {\rm source\ effective\ depuis}\ \delta W_{\rm loc}
& \rho_{\rm Spearman}=0.741 \\
{\rm Paper\ 15,\ Step\ E}
& {\rm source\ radiale\ contrôlée}
& {\rm géométries\ contrôlées}
& \rho_{\rm Spearman}=0.752,\ {\rm near/far}=11.51 \\
\bottomrule
\end{array}
\]

Ainsi, le même couplage
\[
\delta W_{\rm loc}
\longrightarrow
\delta\kappa
\]
apparaît sous trois angles indépendants : réponse locale directe
(Paper~7), reconstruction d'une source effective \(T_{\mu\nu}^{\rm eff}\)
(Paper~8), et réponse radiale contrôlée sur géométries connues
(Paper~15). Le résultat de Paper~15 est particulièrement important car il
confirme ce couplage dans un cadre où la géométrie de référence est contrôlée
et où la courbure attendue est connue.

\section{Assemblage vers l'équation d'Einstein effective}

Si les limites précédentes sont satisfaites, l'action discrète BuP admet une
limite continue schématique :
\[
S_{\rm cont}[g]
=
\int d^Dx\sqrt{|g|}
\left[
\frac{1}{16\pi G_{\rm eff}}R
+
\Lambda_{\rm ent}
+
\mathcal{L}_{\rm ent-source}
+
\text{termes de courbure supérieure}
\right].
\]

La variation par rapport à \(g^{\mu\nu}\) donne :
\[
G_{\mu\nu}
+
\Lambda_{\rm ent}g_{\mu\nu}
+
\mathcal{H}_{\mu\nu}
=
8\pi G_{\rm eff}T_{\mu\nu}^{\rm ent},
\]
où \(\mathcal{H}_{\mu\nu}\) contient les corrections non locales ou de
courbure supérieure.

Dans la limite lisse de basse énergie :
\[
\mathcal{H}_{\mu\nu}\to0,
\]
on obtient :
\[
G_{\mu\nu}
+
\Lambda_{\rm ent}g_{\mu\nu}
=
8\pi G_{\rm eff}T_{\mu\nu}^{\rm ent}.
\]

Cette équation n'est pas postulée : elle apparaît comme limite continue de
l'équilibre variationnel du réseau d'intrication.

\section{Résumé des résultats numériques}

Les trois premiers tests numériques de Paper~15 soutiennent les trois flèches
nécessaires à la dérivation effective :

\[
W_{ij}\to L_\epsilon\to\Delta_g,
\qquad
\kappa_{ij}^{\rm OR}\to R_{\mu\nu}u^\mu u^\nu,
\qquad
\delta W_{\rm loc}\to\delta\kappa(r).
\]

\[
\begin{array}{llll}
\toprule
{\rm Étape} & {\rm Quantité} & {\rm Valeur} & {\rm Statut} \\
\midrule
{\rm Step\ B} & d_s({\rm circle}) & 1.0021 & positif \\
{\rm Step\ B} & d_s({\rm grid2d}) & 1.9938 & positif \\
{\rm Step\ B} & d_s({\rm interval}) & 0.9445 & positif \\
{\rm Step\ B} & d_s({\rm sphere}) & 1.8964 & positif \\
{\rm Step\ C} & \Delta\bar\kappa_{\rm sphere-flat} & 0.010525 & positif\ relatif \\
{\rm Step\ E} & {\rm near/far,\ flat\ torus} & 7.75 & positif \\
{\rm Step\ E} & {\rm near/far,\ sphere} & 11.51,\ \rho=0.752 & positif \\
\bottomrule
\end{array}
\]

Ces résultats ne prouvent pas encore les équations d'Einstein continues. Ils
valident numériquement les briques nécessaires à une dérivation effective.

Le Step~E ne constitue donc pas un résultat isolé : il confirme, sur des
géométries contrôlées, le même couplage \(\delta W_{\rm loc}\to\delta\kappa\)
déjà observé dans Paper~7 et Paper~8.

\section{Limites}

Plusieurs limites doivent être clairement énoncées.

Premièrement, la convergence
\[
L_N\to\Delta_g
\]
est numériquement soutenue sur des géométries contrôlées, mais pas encore
prouvée analytiquement.

Deuxièmement, la correspondance
\[
\kappa_{ij}^{\rm OR}\to R_{\mu\nu}u^\mu u^\nu
\]
est pour l'instant qualitative et relative. Le test montre que la sphère
présente un excès de courbure par rapport au tore plat, mais ne démontre pas
une convergence ponctuelle.

Troisièmement, la source \(T_{\mu\nu}^{\rm ent}\) reste à formaliser
complètement. Paper~8 fournit une construction numérique, mais la dérivation
continue doit encore être reliée rigoureusement à la première loi de
l'intrication :
\[
\delta S_A=\delta\langle K_A\rangle.
\]

Quatrièmement, le terme topologique \(S_{\rm topo}[W]\) n'a pas encore de
forme fermée.

Cinquièmement, l'action continue obtenue peut contenir des corrections
supérieures :
\[
\mathcal{H}_{\mu\nu}\neq0.
\]
Il reste à déterminer le régime où ces corrections deviennent négligeables.

\section{Conclusion}

Paper~15 établit la première étape d'une dérivation effective des équations
d'Einstein depuis l'équilibre d'un réseau d'intrication.

À partir de l'action discrète :
\[
S_{\rm BuP}[W]
=
\mathrm{Tr}\,L(W)^{-\beta[W]}
+
\sum_{ij}W_{ij}\kappa_{ij}[W]
+
\lambda\sum_{ij}W_{ij}d_{ij}^{2}
+
S_{\rm topo}[W],
\]
on étudie la limite :
\[
\frac{\delta S_{\rm BuP}}{\delta W_{ij}}=0
\quad
\xrightarrow[N\to\infty]{}
\quad
G_{\mu\nu}
+
\Lambda_{\rm ent}g_{\mu\nu}
=
8\pi G_{\rm eff}T_{\mu\nu}^{\rm ent}.
\]

Les tests numériques montrent que :
\[
L_\epsilon=(D-W)/\epsilon
\]
reconstruit la dimension spectrale de géométries contrôlées, que la courbure
d'Ollivier--Ricci distingue une géométrie positivement courbée d'une
géométrie plate périodique, et qu'une perturbation locale de \(W_{ij}\)
produit une réponse de courbure localisée.

Ainsi, Einstein n'est pas posé comme axiome dans BuP. Il apparaît comme une
cible de limite continue : la gravité est la forme macroscopique de la
condition variationnelle du réseau d'intrication.

\section*{Disponibilité du code et des données}

Les scripts et résultats de Paper~15 sont organisés dans le dossier :
\[
\texttt{papers/paper15\_einstein\_derivation/}.
\]

Les scripts principaux sont :
\[
\texttt{scripts/paper15\_spectral\_convergence\_v2.py},
\]
\[
\texttt{scripts/paper15\_ricci\_convergence\_v2.py},
\]
\[
\texttt{scripts/paper15\_source\_response\_v2.py}.
\]

Les résultats numériques principaux sont stockés dans :
\[
\texttt{results/spectral\_convergence\_v2\_best\_unnormalized/},
\]
\[
\texttt{results/ricci\_convergence\_v2/},
\]
\[
\texttt{results/source\_response\_v2/}.
\]

\appendix

\section{Commentaires sur la trace de chaleur}

La trace de chaleur d'un opérateur elliptique continu admet l'expansion :
\[
\mathrm{Tr}\,e^{-t\Delta_g}
\sim
(4\pi t)^{-D/2}
\sum_{n=0}^{\infty}
a_n(\Delta_g)t^n.
\]

La représentation de Mellin donne :
\[
\mathrm{Tr}\,L^{-\beta}
=
\frac{1}{\Gamma(\beta)}
\int_0^\infty
t^{\beta-1}
\mathrm{Tr}(e^{-tL})\,dt.
\]

Dans la limite continue, les coefficients de Seeley--DeWitt contiennent les
termes géométriques :
\[
a_0\sim \int d^Dx\sqrt{|g|},
\]
et
\[
a_1\ \text{ou}\ a_2
\sim
\int d^Dx\sqrt{|g|}R,
\]
selon la convention de l'opérateur et la dimension. Ce développement suggère
que le terme spectral BuP contient une contribution de type
Einstein--Hilbert, ainsi que des corrections supérieures.

\section{Prochaines étapes}

Les prochaines étapes sont :

\begin{enumerate}
    \item tester la convergence spectrale avec des tailles plus grandes
    \(N=2048,4096\) en utilisant des méthodes sparse ;
    \item démontrer ou valider numériquement la convergence ponctuelle de
    \(\kappa_{ij}^{\rm OR}\) vers \(R_{\mu\nu}u^\mu u^\nu\) ;
    \item reconstruire explicitement \(g_{\mu\nu}^{\rm ent}\) à partir de
    \(d_{\rm ent}(i,j)\) ;
    \item formaliser \(T_{\mu\nu}^{\rm ent}\) depuis la première loi de
    l'intrication ;
    \item déterminer la forme minimale de \(S_{\rm topo}[W]\) ;
    \item identifier les corrections \(\mathcal{H}_{\mu\nu}\) et leur régime
    de suppression.
\end{enumerate}

\end{document}
